\documentclass[11pt]{article}
\usepackage[T1]{fontenc}
\usepackage{lmodern}
\usepackage[margin=1in]{geometry}
\usepackage{amsmath,amssymb}
\usepackage{microtype}
\usepackage[hidelinks]{hyperref}
\hypersetup{pdftitle={Modulation-speed workload counterexamples for Cox/D/1 queues},pdfauthor={Tynan Daly}}
\newcommand{\E}{\mathbb{E}}
\newcommand{\eps}{\epsilon}
\title{Modulation-speed workload counterexamples\\for Cox/D/1 queues}
\author{Tynan Daly}
\date{10 September 2026}
\begin{document}
\maketitle

\noindent\textbf{Status.} Mathematical derivation with exact-arithmetic certificates and numerical cross-checks, not independent peer review or formal verification. Novelty is not established. The target is the universal workload inequality in Problem~1 of Leskel\"a (2022), including its restriction to finite-state Markov intensity processes. Established positive theorems for special ordered Markov environments are not contradicted.

\section{Main statement}

There is a 64-state continuous-time Markov intensity process with strictly positive intensity levels and strictly positive transition rates between every pair of distinct states such that increasing the modulation speed from $2\pi$ to $3\pi$ strictly increases stationary mean workload in the unit-service, unit-speed single-server queue. An explicit bound is
\[
w(3\pi)-w(2\pi)>\frac{1587}{25\,580\,800\,000\,000}>0.
\]
The environment is stationary, irreducible and mixing. The statement does not rely on a periodic sample path, non-Markov hidden phases, zero intensities or zero transition rates. A separate 24-state example below has entirely rational parameters, though a smaller certified effect; its transition graph is a directed cycle. Neither state count is claimed globally minimal.

\section{The workload inequality used in both certificates}

Let $X$ be stationary with $\E X(0)=1$ and $0\le X\le M$. Conditional on $X$, let arrivals be Poisson of intensity $\eps X(ct)$. Every service requirement is exactly~1, the server speed is~1, and $\eps M<1$. Define
\[
K(c)=\int_0^1(1-u)\E[X(0)X(cu)]\,du.
\]
The stationary mean workload satisfies
\begin{equation}
\frac{\eps}{2}+\eps^2K(c)
\le w(c)\le
\frac{\eps}{2}+\frac{\eps^2K(c)}{1-\eps M}.
\label{eq:bound}
\end{equation}

\subsection*{Construction and integrability}

Use a two-sided stationary environment and an independent homogeneous Poisson proposal process of rate $\eps M$ with independent uniform marks. Accept a point at $t$ when its mark is at most $X(ct)/M$. Start empty at time $-T$. The reflected workloads at zero increase with $T$ and are dominated by the stationary homogeneous M/D/1 workload driven by all proposal points. The latter queue is stable and has finite moments of every order. The monotone limit therefore exists, is stationary by shift-equivariance, and has the required finite moments.

At a deterministic stationary time write $V$ for workload. Integrate the pathwise identity
\[
V(t)^2-V(0)^2=-2\int_0^t V(s)\,ds
+\int_{(0,t]}(2V(s-)+1)N(ds).
\]
Domination and bounded intensity make these terms integrable. The predictable compensator gives
\[
w(c)=\eps/2+\eps\E[X(0)V].
\]
No independence between workload and environment and no PASTA assertion are used.

On the same arrivals define the infinite-server residual work
\[
Y=\sum_{T_i\in(-1,0]}(1+T_i).
\]
Pathwise $Y\le V$: a job of age $u<1$ cannot have received more than $u$ units of service. Conditional Poisson expectation gives $\E Y=\eps/2$ and $\E[X(0)Y]=\eps K(c)$. Scalar stationary autocorrelation is even, so the backwards time argument in this calculation equals the displayed $K(c)$.

Let $D=V-Y\ge0$ and $d=\E D=w(c)-\eps/2$. Then exactly
\[
d=\eps^2K(c)+\eps\E[X(0)D].
\]
The pointwise bounds $0\le X\le M$ imply
\[
\eps^2K(c)\le d\le\eps^2K(c)+\eps M d,
\]
which proves~\eqref{eq:bound}. Thus~\eqref{eq:bound} controls the entire remainder, not merely one term in an asymptotic expansion. In particular,
\begin{equation}
w(c_2)-w(c_1)\ge\eps^2\left[K(c_2)-\frac{K(c_1)}{1-\eps M}\right].
\label{eq:comparison}
\end{equation}

\section{Full-support 64-state Markov generator}

Let $n=64$, $\theta=2\pi/n$, $r=1/\sin\theta$, and $\delta=1/2000$. Let $S$ be the cyclic permutation matrix satisfying $(Sv)_j=v_{j+1\bmod n}$, and let $\Pi$ have every entry equal to $1/n$. Set
\begin{equation}
Q=r(S-I)+\delta(\Pi-I).
\label{eq:generator}
\end{equation}
One interpretation is forward moves around a cycle at rate $r$, with independent uniform redraws of the state at rate $\delta$. A redraw that chooses the same state is a null event. Actual off-diagonal rates are
\[
q_{j,k}=r\,\mathbf{1}_{\{k=j+1\bmod n\}}+\frac{1}{128000},\qquad k\ne j.
\]
The diagonal is minus the off-diagonal row sum. Every off-diagonal entry is strictly positive. Row and column sums vanish, so the stationary distribution is uniform. The finite chain is irreducible and mixing.

Set
\begin{equation}
x_j=1+\frac12\cos\left(\frac{2\pi(j+1/4)}{64}\right),
\qquad X(t)=x_{J(t)}.
\label{eq:levels}
\end{equation}
All 64 $x_j$ are distinct: cosine equality would require either identical indices modulo $n$ or $j+k+1/2$ to be a multiple of $n$, the latter being impossible. Therefore the scalar intensity observation $X(t)$ is itself a CTMC, not merely a non-Markov function of a hidden chain. Its mean is~1 and $1/2<x_j<3/2$.

Use $\eps=1/2000$, $\lambda_c(t)=\eps x_{J(ct)}$, and unit services. The actual environmental generator is $cQ$ and peak arrival intensity is below $3/4000$, so all compared queues are dominated by a stable homogeneous M/D/1 queue.

\subsection*{It lies outside stochastic-monotonicity assumptions}

In the ordering of numerical intensities, $x_{62}<x_1<x_{63}<x_0$. For the upper set $H=\{x_{63},x_0\}$, both states $x_{62}$ and $x_1$ start below $H$. From state~62 the forward move enters $H$, but from state~1 it does not. The uniform-reset contribution to the instantaneous probability of entering $H$ is identical from these two states. Hence
\[
P_{62}(X(t)\in H)-P_1(X(t)\in H)=rt+O(t^2)>0
\]
for sufficiently small $t$. Thus the scalar chain is not stochastically monotone; this is not a counterexample to the established positive theorem under double stochastic monotonicity.

\section{Covariance and exact sign certificate}

On the first Fourier mode, $\Pi$ acts as zero and $S$ acts by $e^{i\theta}$. Its generator eigenvalue is
\[
r(e^{i\theta}-1)-\delta=-\alpha+i,
\qquad \alpha=\tan(\pi/64)+\delta.
\]
Averaging over uniform initial states yields
\[
\E[X(0)X(t)]=1+\frac18e^{-\alpha t}\cos t,\qquad t\ge0.
\]
Consequently
\begin{equation}
K(c)=\frac12+\frac18 I_\alpha(c),\qquad
I_\alpha(c)=\operatorname{Re}\frac{e^z-1-z}{z^2},\qquad z=(-\alpha+i)c.
\label{eq:kernel}
\end{equation}
For completeness, the following certificate uses only elementary inequalities and rational arithmetic, not a floating-point sign test.

Using $3<\pi<22/7$, $\sin x<x$ and $\cos x>1-x^2/2$, with $u=11/224$,
\[
0<\alpha<\frac{u}{1-u^2/2}+\frac1{2000}
=\frac{9956231}{200462000}<\frac1{20}.
\]
Define $A=\alpha/(1+\alpha^2)$ and $C=(1-\alpha^2)/(1+\alpha^2)^2$. Then $A<\alpha$ and $99/100<C<1$. At multiples of $\pi$,
\[
I_\alpha(k\pi)=\frac{A}{k\pi}
+C\frac{1-(-1)^ke^{-\alpha k\pi}}{(k\pi)^2}.
\]
Using $e^{-x}\ge1-x$ and $1-e^{-x}\le x$,
\[
\begin{split}
I_\alpha(3\pi)-I_\alpha(2\pi)
&> -\frac1{360}
+\frac{99}{100}\left[\frac{2}{9(22/7)^2}-\frac1{72}\right]\\
&=\frac{91}{15840}>\frac1{200}.
\end{split}
\]
The bracket is positive, which justifies multiplying its lower bound by the lower bound on $C$. Therefore
\[
K(3\pi)-K(2\pi)>\frac1{1600}.
\]
Also $I_\alpha(2\pi)\le\alpha/\pi<1/60$, giving $K(2\pi)<241/480$. Substitution into~\eqref{eq:comparison} gives
\[
\begin{split}
w(3\pi)-w(2\pi)
&>\frac1{2000^2}
\left[\frac1{1600}-\frac{(3/4000)(241/480)}{1-3/4000}\right]\\
&=\frac{1587}{25\,580\,800\,000\,000}>0.
\end{split}
\]
For orientation, direct formula evaluations give $K(2\pi)$ approximately $0.50182682190600814$ and $K(3\pi)$ approximately $0.50292850997295202$. These decimals are not the sign certificate.

\subsection*{Robustness}

The certificate is strict. In a neighborhood of this full-support generator the stationary distribution, finite-time matrix exponentials and $K$-values depend continuously on the parameters. Intensity labels remain positive and distinct, and the subcritical peak bound remains valid. Normalize perturbed intensity vectors by their stationary mean if keeping $\E X=1$ fixed. The strict separation in~\eqref{eq:comparison} therefore persists in a sufficiently small open neighborhood. No size of that neighborhood is asserted here.

The deterministic service assumption is allowed by the stated problem. This note does not establish a reversal for exponential service times.

\section{A smaller all-rational 24-state example}

Let $Q=S-I$ for a directed 24-cycle with unit forward rates, and use uniform stationarity. Define $x$ by dividing the following exact integers by 1000:
\[
\begin{gathered}
1499,\ 1473,\ 1416,\ 1330,\ 1221,\ 1098,\ 967,\ 839,\\
722,\ 624,\ 552,\ 510,\ 501,\ 527,\ 584,\ 670,\\
779,\ 902,\ 1033,\ 1161,\ 1278,\ 1376,\ 1448,\ 1490.
\end{gathered}
\]
These values define the model; they are not uncertain decimal approximations. They are distinct, positive, average exactly~1 and are below $3/2$. Thus $x_{J(t)}$ itself is an irreducible mixing CTMC. Take $\eps=1/1000000$, unit services, and compare the actual environmental generators $29Q$ and $30Q$.

The verifier obtains rational enclosures of
\begin{equation}
K(c)=\frac1{24}x^T\left[\sum_{m=0}^{\infty}\frac{c^mQ^m}{(m+2)!}\right]x.
\label{eq:series}
\end{equation}
This follows by termwise integration of the uniformly convergent matrix exponential on $[0,1]$. With terms through $m=L$ and $\lVert Q\rVert_\infty=2$, the error in $K$ is at most
\begin{equation}
\frac{M(2c)^{L+1}}{(L+3)!}\frac1{1-2c/(L+4)}
\label{eq:tail}
\end{equation}
when $2c<L+4$. Here $\lVert(x/24)^T\rVert_1=1$ and $\lVert x\rVert_\infty\le M$, which give the scalar norm bound. The ratio of successive omitted terms is at most $2c/(L+4)$, which gives the geometric tail bound.

With $L=240$, exact \texttt{Fraction} arithmetic proves the outward-rounded bounds
\[
K(29)<\frac{503796}{1000000},\qquad
K(30)>\frac{503798}{1000000}.
\]
The omitted-series widths are less than $7\times10^{-51}$ and $3\times10^{-47}$ respectively, but decimal error estimates are not used for the certificate. Applying~\eqref{eq:comparison},
\[
w(30)-w(29)>
\frac{1244303}{999998500000000000000000}
>\frac1{10^{18}}>0.
\]
This is a much smaller absolute certified effect than in the 64-state example. The gain here is fewer states and entirely rational input data, not practical performance magnitude. The smaller state count was found by searching a family of cycle-based examples. Neither failure to find smaller examples nor numerical exploration establishes global minimality.

\section{Reproduction and status}

From the repository root, run
\begin{verbatim}
python3 verify/run_all.py
\end{verbatim}
The scripts use Python 3.10 or newer and its standard library. The runner checks the periodic certificate in Appendix~\ref{sec:periodic}, the 64-state elementary rational certificate, and the 24-state rational matrix-series enclosures. Decimal displays are diagnostics; all comparisons establishing the certificates use exact arithmetic. The final output on success is \texttt{ALL CHECKS PASSED}.

The supplied audit also reports numerical cross-checks of the 64-state generator and its semigroup using NumPy/SciPy, a closed-form integral check with mpmath at 70 decimal digits, and direct matrix-exponential quadrature for the 24-state generator. Those optional numerical routines are not included in this repository. Such numerical checks are separate checks of calculations, not independent review of the probability proof.

\section{Research context and limitations}

Leskel\"a (2022), Problem~1, asks for universal modulation-speed workload monotonicity for stationary ergodic intensities and for a characterization of when it holds. Its discussion records positive results for special finite-state CTMCs, including those whose generator and time reversal are both stochastically monotone. These examples address the universal assertion even within finite-state CTMC environments; they do not contradict those restricted positive theorems.

The simplified random-phase sinusoid is non-Markov and nonmixing. The cycle construction in \path{docs/unrestricted_modulation_speed_v2.md} already avoids those limitations. The full-support refinement above adds positive rates between every pair of distinct states; the 24-state construction reduces state count with an exact computational certificate.

The first Ross conjecture, concerning comparison with a constant-rate Poisson queue, is a different assertion. A nonmonotone dependence on speed does not contradict that lower-bound theorem.

A targeted search and inspection of Leskel\"a's precise statement do not establish originality. Earlier periodic-input results and earlier Markov-modulated queue results remain relevant prior art. Even an old periodic counterexample can be important prior art for a Markov approximation construction. No claim of first discovery, global minimality, full characterization, or all-load reversal is made. Novelty relative to Rolski (1987, 1989), Lemoine (1989), Heyman (1982), and Miyoshi--Rolski (2004) is not established.

\appendix
\section{Random-phase periodic certificate}
\label{sec:periodic}

The following specialization is recorded in \path{docs/upper_bound_audit.md}. Let
\[
X(t)=1+\cos(t+\Theta),\qquad
\Theta\sim\operatorname{Uniform}[0,2\pi),\qquad M=2.
\]
Uniform phase gives stationarity. The continuous-time shifts act as every rotation of the phase circle, so any event invariant under all shifts has probability zero or one: the intensity is ergodic. It is not mixing. The service law is degenerate iid and independent of arrivals.

Direct integration gives
\[
\E[X(0)X(t)]=1+\tfrac12\cos t,\qquad
K(c)=\tfrac12+\frac{1-\cos c}{2c^2}.
\]
At $c_1=2\pi$ and $c_2=3\pi$,
\[
K(c_1)=\tfrac12,\qquad
K(c_2)=\tfrac12+\frac1{9\pi^2}>\tfrac12+\frac1{90}.
\]
The strict elementary estimate uses $\pi^2<10$, for example from $\pi<22/7$. Writing $\rho$ for the mean arrival rate, \eqref{eq:bound} yields
\[
\begin{aligned}
w_\rho(3\pi)-w_\rho(2\pi)
&>\rho^2\left(\frac1{90}-\frac{\rho}{1-2\rho}\right)\\
&=\frac{\rho^2(1-92\rho)}{90(1-2\rho)}.
\end{aligned}
\]
In particular, every $0<\rho<1/92$ is certified by this coarse bound. This is a sufficient range, not a necessary range for the reversal itself.

At $\rho=1/100$,
\[
w_{1/100}(2\pi)\le\frac{99}{19600},\qquad
w_{1/100}(3\pi)>\frac{2273}{450000},
\]
so
\[
w_{1/100}(3\pi)-w_{1/100}(2\pi)>\frac1{11025000}>0.
\]
At $\rho=1/1000$,
\[
w_{1/1000}(2\pi)\le\frac{999}{1996000},\qquad
w_{1/1000}(3\pi)>\frac{22523}{45000000},
\]
so
\[
w_{1/1000}(3\pi)-w_{1/1000}(2\pi)>\frac{227}{22455000000}>0.
\]

\section*{Acknowledgment of AI assistance}

AI tools assisted with the mathematical derivation, writing, and verification code; the author is responsible for the content.

\begin{thebibliography}{9}
\bibitem{leskela}
L. Leskel\"a (2022).
\emph{Ross's second conjecture and supermodular stochastic ordering}.
Queueing Systems 100, 213--215.
\href{https://doi.org/10.1007/s11134-022-09824-0}{doi:10.1007/s11134-022-09824-0};
\href{https://arxiv.org/abs/2204.11856}{arXiv:2204.11856}.
The supplied audit reports direct inspection of the problem statement and CTMC discussion.

\bibitem{baeuerle}
N. Baeuerle and T. Rolski (1998).
\emph{A monotonicity result for the workload in Markov-modulated queues}.
Journal of Applied Probability 35(3), 741--747.
\href{https://doi.org/10.1239/jap/1032265221}{doi:10.1239/jap/1032265221}.
The positive theorem's hypothesis is discussed in reference~\cite{leskela}; this audit does not claim a fresh full-text review of reference~\cite{baeuerle}.

\bibitem{miyoshi}
N. Miyoshi and T. Rolski (2004).
\emph{Ross-type conjectures on monotonicity of queues}.
Australian \& New Zealand Journal of Statistics 46(1), 121--131.
Contains positive comparisons under regularity assumptions; not evidence that every finite-state Markov environment has the ordering.
\end{thebibliography}
\end{document}
